\section{Node.js}
\label{section:Node.js}
Node.js \cite{Node.js} to wieloplatoformowe środowisko uruchomieniowe JavaScript, o otwartym kodzie źródłowym, które umożliwia wykonywanie kodu JavaScript poza przeglądarką. Jest oparte na silniku V8, rozwijanym przez Google, który kompiluje JavaScript do natywnego kodu maszynowego, zapewniając wysoką wydajność.
Node.js jest szeroko używany w aplikacjach serwerowych dzięki swojej asynchroniczności i popularności.  W dodatku JavaScript jest jednym z najbardziej popularnych języków.
\subsection{Działanie Node.js}
\begin{enumerate}
    \item \textbf{Silnik V8:} Node.js korzysta z silnika JavaScript V8, który wykonuje kod JavaScript w szybki i zoptymalizowany sposób, kompilując go do natywnego kodu maszynowego.
    \item \textbf{Model zdarzeniowy:} Node.js używa modelu opartego na zdarzeniach (event-driven model), który pozwala na obsługę wielu żądań jednocześnie bez konieczności blokowania wątków.
    \item \textbf{Asynchroniczność:} Node.js stosuje asynchroniczne wejście/wyjście, co oznacza, że operacje takie jak odczyt plików, zapytania do bazy danych czy żądania sieciowe są wykonywane w tle, bez zatrzymywania głównego wątku.
    \item \textbf{Jednowątkowość:} Mimo że Node.js działa na pojedynczym wątku, korzysta z puli wątków w tle (thread pool) do obsługi operacji wymagających więcej czasu, takich jak operacje na plikach czy połączenia sieciowe.
\end{enumerate}
\subsection{Zalety Node.js}
\begin{enumerate}
    \item \textbf{Wydajność:} Dzięki silnikowi V8 i modelowi asynchronicznemu, Node.js oferuje wysoką wydajność i szybkie przetwarzanie żądań.
    \item \textbf{Jednolity język:} Dzięki Node.js możliwe jest używanie JavaScript po stronie klienta, jak i po stronie serwera. Dzięki temu nie trzeba używać osobnych języków do frontendu jak i backendu.
    \item \textbf{Bogaty ekosystem:} Node.js posiada ogromną liczbę modułów i bibliotek dostępnych w rejestrze npm (Node Package Manager), co pozwala na szybkie i łatwe rozszerzanie funkcjonalności aplikacji. Jedną z nich jest net jak i ws. Net pozwala na połączenia przez  protokół TCP z ip w sieci, natomiast ws pozwala tworzyć WebSockety. Biblioteka net została użyta do komunikacji z serwerem na STM32, ws zostało zastosowane do komunikacji z aplikacją webową na frontend.
    \item \textbf{Asynchroniczność:} Asynchroniczne wejście/wyjście sprawia, że aplikacje oparte na Node.js są bardziej responsywne i skalowalne.
    \item \textbf{Społeczność:} Duża społeczność i szerokie wsparcie umożliwiają szybkie rozwiązywanie problemów i korzystanie z gotowych rozwiązań jakimi są biblioteki czy kody do obsługi ów bibliotek.
    \item \textbf{Wsparcie dla mikroserwisów:} Node.js dobrze nadaje się do tworzenia nowoczesnych aplikacji opartych na mikroserwisach. Mikroserwisy zostały w swoisty sposób użyte w projekcie, pobieranie próbek i obsługa rotatorów na STM32, pośrednik z WebSocketem na Node.js do komunikacji pomiędzy STM32 i aplikacją webową i aplikacja webowa jako interfejs użytkownika do obsługi zarówno rotatorów jak i analizie danej pozycji.
\end{enumerate}
\subsection{Podsumowanie}
Node.js jest powszechnie stosowanym środowiskiem uruchomieniowym JavaScript, które dzięki swojej wydajności, asynchroniczności i popularności znajduje szerokie zastosowanie w aplikacjach serwerowych, aplikacjach czasu rzeczywistego, systemach IoT jak i w wielu innych zastosowaniach. Jego zaletą jest wysoka wydajność i bogaty ekosystem.

\section{NPM}
\label{section:NPM}
NPM (Node Package Manager) to domyślny menedżer pakietów dostarczanych wraz z Node.js. Jest również największym rejestrem bibliotek o otwartym kodzie źródłowym na świecie, które są głównie wykorzystywanie do zarządzania zależnościami w projektach opartych zarówno na JavaScript jak i na Node.js. Umożliwia instalowanie, zarządzanie jak i aktualizowanie bibliotek i modułów w prosty i wydajny sposób.
\subsection{Działanie npm}
npm działa głównie w trzech obszarach:
\begin{enumerate}
    \item \textbf{Rejestr pakietów:} Centralne repozytorium, w którym przechowywane są moduły i biblioteki do zarówno JavaScript jak i Node.js. Daje to możliwość publikowania własnych modułów jak i korzystania z istniejących pakietów.
    \item \textbf{Klient wiersza poleceń:} Umożliwia użytkownikom interakcję z rejestrem za pomocą poleceń, takich jak instalowanie, aktualizowanie, usuwanie i zarządzanie pakietami.
    Umożliwia użytkownikom komunikację z rejestrem za pomocą poleń, takich jak instalowanie, aktualizowanie, usuwanie i zarządzanie pakietami. Komunikacja jest dostępna za pomocą komend w wierszu poleceń takich jak cmd, powershell czy linux shell.
    \item \textbf{Plik package.json}: Plik konfiguracyjny w formacie JSON, w którym przechowywane sa informacje o zależnościach projektu jak i skryptach uruchamianych przez npm.
\end{enumerate}
Oprócz package.json istnieje też plik package-lock.json, który umożliwia:
\begin{enumerate}
    \item \textbf{Zablokowanie wersji pakietów:} Plik package-lock.json przechowuje w formacie JSON szczegółowe dane o wersjach wszystkich zainstalowanych pakietów i ich zależności, nawet jeśli w package.json nie była podana wersja lub został podany zakres wersji. Dzięki temu przy instalacji zależności aplikacja będzie zawsze miała te same wersje pakietów, co pozwala uniknąć potencjalnych problemów wynikających z aktualizacji pakietów.
    \item \textbf{Stabilność projektu:}  Podczas instalacji pakietów i zależności w pierwszej kolejności są wykorzystywanie informacje zawarte w package-lock.json, co gwarantuje, że środowisko projektowe będzie korzystało z tych samych pakietów na różnych urządzeniach niezależnie od czasu instalacji. Jest to kluczowe przy instalacji aplikacji na wielu urządzeniach. Nawet drobne zmiany w wersjach paczek i zależności mogą prowadzić do nieprzewidywalnych problemów z działaniem aplikacji.
    \item \textbf{Poprawa wydajności:} Dzięki package-lock.json npm może szybciej instalować pakiety, ponieważ dokładnie wie, które wersje należy pobrać, zamiast analizować i rozwiązywać zależności od nowa.
\end{enumerate}
\subsection{Zalety \texttt{npm}}
\begin{enumerate}
    \item \textbf{Obszerny ekosystem:} npm posiada niezwykle dużą bazę bibliotek, co pozwala na znalezienie sprawdzonego i gotowego rozwiązania, zamiast samemu pisać wszystko od zera.
    \item \textbf{Łatwość użycia:} łatwe i intuicyjne komendy w wierszu poleceń sprawiają, że używanie obsługa npm jest bardzo prosta.
    \item \textbf{Zarządzanie wersjami:} Umożliwia precyzyjne kontrolowanie wersji pakietów dzięki semantycznemu wersjonowaniu (SemVer).
    \item \textbf{Integracja z Node.js:} npm jest domyślnie zintegrowany z Node.js, co czyni go standardowym narzędziem w ekosystemie JavaScript.
    \item \textbf{Automatyzacja:} dzięki skryptom w package.json, npm pozwala na pełną automatyzację zadań takich jak uruchamianie aplikacji jak i produkcyjne budowanie aplikacji.
    \item \textbf{Wsparcie społeczności:} Duża społeczność użytkowników i aktywnych deweloperów zapewnia szybkie wsparcie oraz regularne aktualizacje.
    Ogromna społeczność, która codziennie zapewnia wsparcie jak i regularnie aktualizuje paczki.
\end{enumerate}
\subsection{Komendy npm}
Komendy dostępne do użycia w aplikacji webowej.
\begin{enumerate}
    \item \textbf{npm start:}  
    Uruchamia aplikację w trybie deweloperskim (development mode).  
    \begin{itemize}
        \item Serwer deweloperski na http://localhost:3000 zostaje uruchomiony, a aplikacja automatycznie otwiera się w przeglądarce. W przypadku gdy port 3000 jest zajęty aplikacja uruchomi się na następnym wolnym porcie np. 3001.
        \item Funkcja Hot Module Replacement (HMR) zapewnia automatyczne odświeżanie strony w przypadku zmian w kodzie, co przyspiesza proces tworzenia aplikacji jak i testowaniu małych zmian np. poprawieniu literówki.
        \item Nadaje się do pracy nad aplikacją lokalnie, ale nie jest przeznaczony do wdrażania w środowisku produkcyjnym.
    \end{itemize}
    \item \textbf{npm test:}  
    Uruchamia zestaw testów aplikacji przy użyciu frameworka \textit{Jest}. Ta komenda nie była używana w czasie prac nad aplikacją webową.
    \item \textbf{npm run build:}  
    Generuje zoptymalizowaną wersję aplikacji przeznaczoną do środowiska produkcyjnego. Funkcjonalność tej komendy zostanie opisana w dalszej części pracy.
\end{enumerate}
\subsection{Komenda npm run build}
Jest to kluczowa komenda w procesie przygotowania aplikacji do wdrożenia do środowiska produkcyjnego lub wydania zoptymalizowanej wersji. Oto co się dzieje podczas jej uruchomienia:
\begin{enumerate}
    \item \textbf{Kompilacja kodu:}
    Komenda przekształca kod źródłowy React, który jest napisany w JSX, na zoptymalizowany kod JavaScript, który jest możliwy do uruchomienia w przeglądarkach bez wcześniejszej interpretacji.
    
    \item \textbf{Zoptymalizowanie zasobów:}  
    Pliki statyczne, takie jak obrazy czy pliki stylów CSS, są minimalizowane w celu zmniejszenia ich rozmiarów. Kod JavaScript również jest minimalizowany i rozdzielany na mniejsze części (chunks) co pozwala na ładowanie tylko niezbędnych modułów zamiast wszystkich plików skryptowych.

    \item \textbf{Generowanie plików wynikowych:}  
    Wszystkie pliki wynikowe są umieszczane w folderze \texttt{build}. Folder ten zawiera:
    Wszystkie wygenerowane pliki są umieszczane w folderze build. Ów folder zawiera między innymi:
    \begin{itemize}
        \item \textbf{index.html} – główny plik HTML, który ładuje aplikację.
        \item \textbf{static/} – folder z minimalizowanymi zasobami (JavaScript, CSS, zdjęcia itp.).
    \end{itemize}
\end{enumerate}
\subsection{Dodatkowe informacje zbudowanej aplikacji}
Po wygenerowaniu zoptymalizowanej wersji aplikacji nie można w standardowy sposób skorzystać z pliku index.html do załadowania aplikacji.

\subsection*{Ograniczenia w działaniu}
Plik index.html odwołuje się do ścieżek z zasobami JavaScript i CSS za pomocą względnych ścieżek, które wymagają serwera HTTP do poprawnego działania. Otwierając bezpośrednio plik index.html przeglądarka nie jest w stanie poprawnie załadować zasobów przez względne ścieżki.

\subsection*{Uruchomienie zbudowanej aplikacji}
Aby poprawnie uruchomić zbudowaną aplikację, należy użyć serwera HTTP. Najprostszym rozwiązaniem jest użycie komendy npx serve build, który uruchomi serwer na podstawie plików w folderze build. npx jest narzędziem do uruchamiania pakietów Node.js, które jest dołączone do npm. Jego głównym celem jest umożliwienie łatwego uruchamiania programów jak i narzędzi, które są zainstalowane w projekcie. npx został również użyty do zbudowania szkieletu aplikacji. Innym rozwiązaniem jest napisanie serwera na Node.js, jednak jest to o wiele bardziej czasochłonne rozwiązanie pod względem wdrożenia. W projekcie został użyty pakiet \textit{serve}, w celu włączenia zbudowanej aplikacji należy użyć polecenia \textit{serve -s build}.